\documentclass[10pt,a4paper]{amsart}
\usepackage{amsmath,amssymb,amsthm,graphicx}
\usepackage{hyperref}
\numberwithin{equation}{section}

%\usepackage{subfigure}
\usepackage{upgreek}
\usepackage{mathtools}
\usepackage{tikz}
\usepackage[all]{xy}
\usetikzlibrary{matrix,calc,arrows.meta,positioning,patterns}


\newcommand{\bC}{{\mathbb C}}
\newcommand{\bP}{{\mathbb P}}
\newcommand{\bQ}{{\mathbb Q}}
\newcommand{\bR}{{\mathbb R}}
\newcommand{\bZ}{{\mathbb Z}}
\newcommand{\bN}{{\mathbb N}}
\newcommand{\bS}{{\mathbb S}}
\newcommand{\bF}{{\mathbb F}}
\newcommand{\bT}{{\mathbb T}}
\newcommand{\bK}{{\mathbb K}}
\newcommand{\bfk}{{\mathbf k}}
\newcommand{\bfq}{{\mathbf q}}


\newcommand{\cF}{\mathcal F}
\newcommand{\cG}{\mathcal G}
\newcommand{\cA}{\mathcal A}
\newcommand{\cB}{\mathcal B}
\newcommand{\cC}{\mathcal C}
\newcommand{\cD}{\mathcal D}
\newcommand{\cE}{\mathcal E}
\newcommand{\cK}{\mathcal{K}}
\newcommand{\cY}{\mathcal Y}
\newcommand{\cS}{\mathcal S}
\newcommand{\cN}{\mathcal N}
\newcommand{\cM}{\mathcal M}
\newcommand{\cO}{\mathcal O}
\newcommand{\cP}{\mathcal P}
\newcommand{\cU}{\mathcal U}
\newcommand{\cL}{\mathcal L}
\newcommand{\cR}{\mathcal R}
\newcommand{\cT}{\mathcal T}
\newcommand{\cH}{\mathcal H}
\newcommand{\oD}{\overline{D}}
\newcommand{\scrS}{\EuScript S}
\newcommand{\scrC}{\EuScript C}
\newcommand{\scrD}{\EuScript D}
\newcommand{\scrE}{\EuScript E}
\newcommand{\scrF}{\EuScript F}
\newcommand{\scrG}{\EuScript G}
\newcommand{\scrL}{\EuScript L}
\newcommand{\scrM}{\EuScript M}
\newcommand{\scrO}{\EuScript O}
\newcommand{\scrP}{\EuScript P}
\newcommand{\scrU}{\EuScript U}
\newcommand\isomto{\stackrel{\sim}{\smash{\longrightarrow}\rule{0pt}{0.4ex}}}
\newcommand{\ev}{\operatorname{ev}}
\DeclareMathOperator{\val}{val}
\DeclareMathOperator{\skel}{Skel}
\DeclareMathOperator{\id}{id}
\DeclareMathOperator{\im}{im}
\DeclareMathOperator{\hoSC}{hoSC}
\DeclareMathOperator{\whoSC}{\widehat{hoSC}}
\DeclareMathOperator{\SC}{SC}
\DeclareMathOperator{\wSC}{\widehat{SC}}
\DeclareMathOperator{\wtel}{\widehat{tel}}
\DeclareMathOperator{\Sp}{span}
\DeclareMathOperator{\tel}{tel}
\DeclareMathOperator{\inj}{inj}
\DeclareMathOperator{\conv}{conv}
\DeclareMathOperator{\Sec}{sec}
\DeclareMathOperator{\Area}{Area}
\DeclareMathOperator{\Vol}{Vol}
\DeclareMathOperator{\diam}{diam}
\DeclareMathOperator{\Tot}{Tot}
\newtheorem{Theorem}{Theorem}[section]
\newtheorem*{Theorem*}{Theorem}
\newtheorem{Corollary}[Theorem]{Corollary}
\newtheorem{Lemma}[Theorem]{Lemma}
\newtheorem{tm}[Theorem]{Theorem}
\newtheorem{lm}[Theorem]{Lemma}
\newtheorem{prp}[Theorem]{Proposition}
\newtheorem{Proposition}[Theorem]{Proposition}
\newtheorem{cy}[Theorem]{Corollary}

\theoremstyle{definition}
\newtheorem{Definition}[Theorem]{Definition}
\newtheorem{Note}[Theorem]{Note}
\newtheorem{Example}[Theorem]{Example}
\newtheorem{Remark}[Theorem]{Remark}
\newtheorem{df}[Theorem]{Definition}
\newtheorem{ques}[Theorem]{Question}
\newtheorem{ass}[Theorem]{Assumption}
\newtheorem{obs}[Theorem]{Observation}
\newtheorem{ex}[Theorem]{Example}
\newtheorem{rem}[Theorem]{Remark}
\usepackage{cleveref}
\crefname{originalreference}{original reference}{original references}
\Crefname{originalreference}{Original reference}{Original references}
\makeatletter
\newlabel{prop:starshape-hypersurface}{{1.22}{source}{}{sourcepointer-1}{}}
\newlabel{prop:starshape-hypersurface@cref}{{[originalreference][0][]1.22}{[1][0][]source}{}{}{}}
\makeatother
\title{Quantitative toric flux-distance rigidity and geometric boundedness of convex torus domains}
\author{Yoel Groman}
\date{doi:10.3842/SIGMA.2026.066; CC BY 4.0}
\begin{document}
\maketitle
\noindent\textit{Allowed standard prerequisites (editorial).}\begin{itemize}
\item Symplectic cutting, holomorphic-curve compactness and standard toric/Fukaya background as used by author. Original Sikorav, STVFluxShape and FOOO/ChoOh bibliography retained; actual source-local arguments included.
\end{itemize}\medskip
\ifdefined\hypertarget\hypertarget{sourcepointer-1}{}\fi
\noindent\textit{Source reference.} Original reference 1.22: algebra-geometry/sources/sigma-2026-066/native/sigma26-066.tex, source lines 381--381. Later application or generalization mentioned in the narrative; not an input to the retained target proof.\par\medskip
\subsection{The results}
As our first proxy for completeness, we use the well-established notion of \emph{geometric boundedness} (see, e.g., \cite[Section~4.2]{Sikorav}): an~$\omega$-compatible almost complex structure $J$ on $M$ is called \emph{geometrically bounded} if the associated Riemannian metric $g_J(\cdot,\cdot)=\omega(\cdot,J\cdot)$ is complete and has uniformly bounded sectional curvature and uniformly positive injectivity radius.\footnote{To treat $\omega$-tame almost complex structures, one must also impose a~uniform bound on the norm of $\omega$ with respect to $g_J$ (i.e., $|\omega|_{g_J}\le C$).}

We begin with an existence question.
\begin{ques}\label{ques:1}
Given $\omega$, does there exist a~geometrically bounded $\omega$-compatible almost complex structure $J$ on $M$?
\end{ques}
It is customary to call $(M,\omega)$ itself geometrically bounded if the answer to Question~\ref{ques:1} is affirmative. We consider this question when we start with a~geometrically bounded $M$ and excise a~subset $N\subset M$. Question~\ref{ques:1} can be understood as the question: \emph{can the distance to $N$ of any compact subset of $M\setminus N$ be made arbitrarily large by adjusting the almost complex structure?}



\begin{tm}\label{tm:1}
Let $M$ be a~geometrically bounded symplectic manifold with finitely generated spherical homology, and suppose $\omega$ is rational.
Let $N\subset M$ be a~properly embedded symplectic hypersurface.
Then $M\setminus N$ is not geometrically bounded.
\end{tm}

\begin{tm}\label{tm:2}
Let $U\subset \bR^n$ be convex.
Let $\omega$ be the standard symplectic form on $\bR^n\times\bT^n$. If there exists an $\omega$-compatible almost complex structure on $U\times\bT^n$ which is geometrically bounded, then $U=\bR^n$.
\end{tm}

Theorem~\ref{tm:2} follows from the following more quantitative statement which is of independent interest.
Identify primitive classes in $H_1(\bT^n;\bZ)$ with primitive vectors $v\in\bZ^n$, hence with integral conormals.
For an open convex set $U\subset\bR^n$ and $p\in U$, let $L_p := \{p\}\times \bT^n \subset U\times\bT^n$. Define the ``flux distance'' to the boundary by
\[
\delta_{\mathrm{flux}}(p;\partial U):=\inf_{v\in\bZ^n \mathrm{primitive}}\sup_{x\in \partial U}\langle v,p-x\rangle\in(0,+\infty).
\]

\begin{Proposition}\label{prop:flux-distance}
Fix $a>0$. There exists a~universal dimensionless constant $C>0$ such that for every open convex $U\subset\bR^n$, every $p\in U$, and every $\omega$-compatible almost complex structure $J$ on $\bR^n\times\bT^n$ for the standard symplectic form with $g_J$ $a$-bounded $($see Definition~{\rm\ref{df:a-bounded}} below$)$, one has
\[
 d_{g_J}(L_p, (\bR^n\setminus U)\times\bT^n)\leq Ca\delta_{\mathrm{flux}}(p;\partial U)
\]
provided $\delta_{\mathrm{flux}}(p;\partial U)\gg 1/a$.
\end{Proposition}
This is proved in Section~\ref{sec:proofs} and immediately implies Theorem~\ref{tm:2}.

\section{Proofs}\label{sec:proofs}

\subsection[Energy confinement for J-holomorphic curves]{Energy confinement for $\boldsymbol{J}$-holomorphic curves}
Throughout, if $u\colon (\Sigma,\partial\Sigma)\to (M,\omega)$ is a~$J$-holomorphic curve with (possibly empty) boundary, its energy is
$
E(u):=\int_{\Sigma} u^{*}\omega \ge 0$.


The following confinement lemmas are stated for $\omega$-compatible $J$. At the end of this subsection we note that both hold for $\omega$-tame $J$ under the uniform tameness condition from the footnote to the definition of geometric boundedness on page \pageref{tm:1} (in Section 1.1).


\begin{lm}[Sikorav's confinement estimate]\label{lem:sikorav}
 Let $(M,\omega)$ be a~symplectic manifold and let $J$ be an $\omega$-compatible almost complex structure such that the associated metric $g_J(\cdot,\cdot):=\omega(\cdot,J\cdot)$ is complete and has uniformly bounded sectional curvature and positive injectivity radius $($i.e., $J$ is geometrically bounded$)$.
 Then there exist a~constant $r_0>0$ depending on the geometry and a~universal constant $c>0$ such that for every connected $J$-holomorphic curve $u$ one has
 \[
 u(\Sigma)\subset B_{r_0\,\lceil E(u)/(c r_0^2)\rceil}(u(\partial\Sigma)),
 \]
 where $B_{r}(\cdot)$ denotes the closed $g_J$-ball of radius $r$ about a~subset and $\lceil\cdot\rceil$ is the ceiling function.
 \end{lm}
 \begin{proof}
 See \cite[Proposition 4.4.1]{Sikorav}.
 \end{proof}
 This confinement estimate is concerned with the large scale. Namely, it contains no information about curves with energy much smaller than $cr_0^2$ other than that the diameter is $\leq r_0$. In some cases, e.g., in the proof of Proposition~\ref{prop:starshape-hypersurface}, we need an estimate for the small scale. This is taken care of by the following lemma.
 \begin{lm}[improved confinement]\label{lem:sikorav-improved}
 In the setting of Lemma~{\rm\ref{lem:sikorav}}, with the same $r_0$ and $c$, and letting $C := 1/c$, one has
 \[
 u(\Sigma)\subset B_{\rho}(u(\partial\Sigma)),\qquad \rho:=\max\bigl(3CE(u)/r_0,\sqrt{CE(u)}\bigr).
 \]
 \end{lm}
 \begin{proof}
 The argument relies on the same monotonicity inequality used for Sikorav's lemma: if a~$J$-holomorphic curve meets the center of a~ball of radius $r\le r_0$ and the boundary of the $J$-holomorphic curve does not meet the interior of the ball, then the energy in that ball is at least $c r^2$.

 If $E(u)/\bigl(c r_0^2\bigr)\le 1$, we apply the monotonicity inequality directly: if the image contained a~point at the center of a~ball of radius greater than \smash{$\sqrt{E(u)/c}$} whose interior did not meet $u(\partial\Sigma)$, then monotonicity would force $E(u)\ge c r^2>E(u)$ (since $E(u)\ge c r^2$ implies $r\le\sqrt{E(u)/c}$). Hence $u(\Sigma)$ is contained in a~ball of radius \smash{$\sqrt{E(u)/c} = \sqrt{C E(u)}$} about $u(\partial\Sigma)$.

 If $E(u)/\bigl(c r_0^2\bigr)>1$, take the maximum number $n$ of disjoint balls of radius $r_0$ that the image of $u$ can meet. Then $n\le E(u)/\bigl(c r_0^2\bigr)$ by monotonicity, and the image is contained in a~ball of radius $(2n+1)r_0\le r_0\bigl(1+2E(u)/\bigl(c r_0^2\bigr)\bigr)<3E(u)/(cr_0) = 3CE(u)/r_0$ around $u(\partial\Sigma)$.
 \end{proof}

For later use, we also record a~localized version, originally observed in~\cite{GromanFloerOpen}.
\begin{df}\label{df:a-bounded}
 For $a>0$, we say that the complete Riemannian metric $g_J$ is \emph{$a$-bounded at a~point} $p\in M$ if \smash{$\inj_{g_J}(p)\ge \frac{1}{a}$} and $\Sec_{g_J}\le a^{2}$ on $B_{1/a}(p)$. We say that $g_J$ is \emph{$a$-bounded on a~set} if it is $a$-bounded at every point of that set.
\end{df}
\begin{rem}
 When $g_J$ is $a$-bounded, we can take $r_0=1/(2a)$ \cite[Theorem 4.4.1 and Remark~4.4.3]{McDuffSalamonJHol}.
\end{rem}

\begin{lm}[local confinement]\label{lem:local-confinement}
Fix $a>0$. There exist $r_0=r_0(a)>0$ and a~universal constant $C>0$ with the following property. Let $K\subset M$ be compact and let $R>0$.
Assume~$J$ is $\omega$-compatible and $g_J$ is complete and is $a$-bounded on the $g_J$-ball $B_{R}(K)$.
Let $u$ be a~$J$-holomorphic curve $u$ with $u(\partial\Sigma)\subset K$ and write
\[
\rho:=\max\bigl(2CE(u)/r_0,\sqrt{CE(u)}\bigr).
\]
If $\rho\le R$, then $u(\Sigma)\subset B_R(K)$.
\end{lm}
\begin{proof}
The proof of Lemma~\ref{lem:sikorav-improved} relies on the monotonicity inequality which depends only on the local geometry of a~complete Riemannian metric. So taking $r_0$ and $C$ as in Lemma~\ref{lem:sikorav-improved}, we obtain by the same argument that
$
u(\Sigma)\subset B_{\rho}(u(\partial\Sigma))
$.
\end{proof}

\subsection{Proof of Theorem~\ref{tm:2}}

Let $X$ be a~closed symplectic toric manifold with moment polytope $P\subset \bR^n$ and moment map~${\mu\colon X\to P}$.
For $p\in \operatorname{Int} P$, let $L_p:=\mu^{-1}(p)\cong \bT^n$. A~\emph{basic $($Maslov index $2)$ class} in~$\pi_2(X,L_p)$ is a~relative class $\beta_F$ corresponding to a~facet $F$ of $P$.
Geometrically, $\beta_F$ can be constructed as follows. Let $v_F \in \mathbb{Z}^n$ be the primitive integral inward-pointing normal vector to~$F$. The circle action generated by $v_F$ acts freely on the interior of $P$ but collapses to a~point over the facet~$F$. Thus, if we draw a~line segment $L$ in the moment polytope $P$ from $p$ to a~point in the interior of $F$ and, transporting along the orbit of this action as one traverses $L$, we obtain a~disc in the class $\beta_F$.


\begin{Proposition}\label{prop:toric-disc-Psi}
For any $p\in \operatorname{Int} P$ and any $\omega$-tame almost complex structure $J$ on $X$, there exists a~stable $J$-holomorphic map with boundary on $L_p$ whose domain has at least one disc component $u_D\colon (D,\partial D)\to (X,L_p)$ of positive symplectic area such that
\[
\int_D u^*\omega \le \min_{F}\omega(\beta_F),
\]
where the minimum is taken over facets $F$ of $P$.
\end{Proposition}
\begin{rem}
 The inequality above is stated for the \emph{disc component} $u_D$: The configuration guaranteed by the proposition may contain sphere bubbles. In that case, the disc component has strictly smaller area than $\omega(\beta_F)$.
 \end{rem}
\begin{proof}
We use the $\Psi$-invariant of a~Lagrangian submanifold considered by Shelukhin--Tonko\-nog--Vianna~\cite{STVFluxShape}.
Fix an $\omega$-tame almost complex structure $J$ and auxiliary choices (virtual perturbations) defining the Fukaya $A_\infty$-algebra of $L_p$, and pass to a~classically minimal model.
The quantity $\Psi(L_p)\in(0,+\infty]$ is defined as an infimum of valuations of the Maurer--Cartan prepotential $m\bigl(e^{b}\bigr)$; it is independent of all choices (in particular of $J$) by~\cite[Theorem~4.2]{STVFluxShape}.
Furthermore, expanding the definition, $\Psi(L_p)$ is the minimum of the symplectic areas $\omega(\beta)$ over those relative classes $\beta\in H_{2}(X,L_p;\bZ)$ which contribute nontrivially to some symmetrized $A_\infty$ operation; see~\cite[Theorem~4.8]{STVFluxShape}. Note this includes the classes which contribute to the curvature term.

For any such class $\beta$ contributing to the curvature term there exists a~stable $J$-holomorphic map $u$ representing $\beta$ with at least one nonconstant disc component.\footnote{To see this, note that otherwise $\beta$ would be a~purely spherical class. The invariance of the indicator function~\cite[Proposition~5.1]{STVFluxShape} would then dictate that it must also realize the $\psi$ invariant for $J_{\mathrm{toric}}$.
However, for $J_{\mathrm{toric}}$, all non-constant holomorphic spheres are confined to the toric boundary divisor and are strictly disjoint from the interior fiber $L_p$, so their contributions to the $A_\infty$-algebra of $L_p$ necessarily vanish.
Hence, $\partial \beta \neq 0$.}
The class area decomposes as a~sum of the symplectic areas of the disc and sphere components of $u$, so every nonconstant disc component $u_D$ satisfies
\[
\int_D u_D^*\omega \le \omega(\beta),
\]
with strict inequality whenever $u$ also has a~nonconstant sphere component.

Now take the standard toric almost complex structure $J_{\mathrm{toric}}$.
Fukaya--Oh--Ohta--Ono~\cite{FOOO_Toric1} (and earlier Cho--Oh~\cite{ChoOhToric}) compute the leading-order terms of the toric superpotential/potential function: for each facet $F$ the basic Maslov index $2$ class $\beta_F$ contributes nontrivially to the curvature term (equivalently, to the potential), with coefficient $1$; see \cite[Section~11]{FOOO_Toric1} and~\cite{ChoOhToric}.
Therefore, for a~facet $F$ minimizing $\omega(\beta_F)$, the class $\beta_F$ witnesses $\Psi(L_p)$ for $J_{\mathrm{toric}}$, and hence
\[
\Psi(L_p) \le \omega(\beta_F) = \min_{F'}\omega(\beta_{F'}).
\]
Since $\Psi(L_p)$ is independent of $J$, the same inequality holds for our given $J$.
Applying \cite[Theorem~4.8]{STVFluxShape} for this $J$, we obtain a~stable $J$-holomorphic map representing some class $\beta$ with~${\omega(\beta)=\Psi(L_p)\le \min_{F'}\omega(\beta_{F'})}$.
Choosing a~disc component $u_D$ of this stable map, the discussion above gives
\[
\int_D u_D^*\omega\le \omega(\beta)=\Psi(L_p)\le \min_{F'}\omega(\beta_{F'}),
\]
which is the claimed inequality.
\end{proof}

To prove Theorem~\ref{tm:2}, it suffices to prove Proposition~\ref{prop:flux-distance} from the introduction.

Write $\pi\colon \bR^{n}\times\bT^{n}\to \bR^{n}$ for the projection, and for $q\in\bR^n$ write $L_q:=\{q\}\times\bT^n\subset \bR^n\times\bT^n$.
Recall the definition of $\delta_{\mathrm{flux}}(p;\partial U)$ from the introduction.
For later use, given $v\in\bZ^n$ we write
\[
c_{\partial U}(v):=\inf_{x\in \partial U}\langle v,x\rangle\in[-\infty,+\infty),
\qquad
\delta_{\partial U}(p;v):=\langle v,p\rangle-c_{\partial U}(v)\in(0,+\infty],
\]
so $\delta_{\partial U}(p;v)<+\infty$ precisely when $c_{\partial U}(v)>-\infty$.

We will deduce Theorem~\ref{tm:2} and Proposition~\ref{prop:flux-distance} from the following stronger version of Proposition~\ref{prop:flux-distance}.

\begin{Proposition}\label{prop:flux-distance-strong}
Let $U\subset\bR^n$ be an open convex set, and let $V$ be an open neighborhood of~${U\times\bT^n}$ in $\bR^n\times\bT^n$. Let $J$ be an $\omega$-tame almost complex structure on $V$ such that the associated metric $g_J$ is complete on $V$ and is $a$-bounded on $U\times\bT^n$ for some $a>0$. Then for every $p\in U$,
\[
 d_{g_J}(L_p, V\setminus (U\times\bT^n)) \le\max\bigl(4Ca\delta_{\mathrm{flux}}(p;\partial U),\sqrt{C\delta_{\mathrm{flux}}(p;\partial U)}\bigr),
\]
where $C$ is a~universal constant.
\end{Proposition}

To prove Proposition~\ref{prop:flux-distance-strong}, we formulate a~basic lemma.
 \begin{lm}\label{lemma2.11}
 Let $U\subset\bR^n$ be an open convex set, let ${K\subset U}$ be compact, and let ${v\in\bZ^n}$ be primitive with $c_{\partial U}(v)>-\infty$.
 Then there exists a~compact Delzant polytope $P$ such that $K\subset \operatorname{Int} P\subset U$ and such that one facet of $P$ has outward primitive normal $v$ {\rm(}equivalently, is parallel to the supporting hyperplane $\{\langle v,\cdot\rangle=c_{\partial U}(v)\}${\rm)}.
 \end{lm}
\begin{proof}
Choose $\varepsilon>0$ such that the closed $\varepsilon$-neighborhood $K_\varepsilon:=\{x\mid d(x,K)\le \varepsilon\}$ is still contained in $U$ (possible since $U$ is open and $K$ is compact).
Pick $c<c_{\partial U}(v)$ close enough to~$c_{\partial U}(v)$ so that $K_\varepsilon\subset\{x\mid \langle v,x\rangle>c\}$, and set $H:=\{x\mid \langle v,x\rangle=c\}$ and $H^-:=\{x\mid \langle v,x\rangle\ge c\}$.

Since $K_\varepsilon$ is compact, we can find finitely many supporting half-spaces with \emph{rational} (hence integral after scaling) outward normals whose intersection contains $K_\varepsilon$ and is contained in $U$. Intersecting these half-spaces together with $H^{-}$ yields a~compact simple rational polytope $P_0$ with $K\subset \operatorname{Int} P_0\subset U$ and with a~facet parallel to $H$.

Finally, by applying arbitrarily small \emph{corner cuts} one can make the polytope smooth/Delzant while keeping the containments $K\subset \operatorname{Int} P\subset U$ and without changing the facet with outward normal $v$.
\end{proof}
\begin{proof}[Proof of Proposition~\ref{prop:flux-distance-strong}]
Let $\delta:=\delta_{\mathrm{flux}}(p;\partial U)$.
Fix $\varepsilon>0$, and choose a~primitive $v\in\bZ^n$ with $c_{\partial U}(v)>-\infty$ and $\delta_{\partial U}(p;v)\le \delta+\varepsilon$.
Put $c:=c_{\partial U}(v)$ and $\ell(x):=\langle v,x\rangle-c$, so $\ell(p)=\delta_{\partial U}(p;v)$.

Let $J$ be as in the statement. We will prove the inequality
\begin{equation}\label{eqPreciseVersion}
 d_{g_J}(L_p, V\setminus (U\times\bT^n)) \le \max\bigl(4Ca\delta_{\mathrm{flux}}(p;\partial U),\sqrt{C\delta_{\mathrm{flux}}(p;\partial U)}\bigr).
\end{equation}
Let $R>0$ be such that the closed metric ball $B:=B_R(L_p)$ with respect to $g_J$ is contained in~${U\times\bT^n}$.
Since $g_J$ is complete on $V$ and $B\subset U\times\bT^n\subset V$, the closed ball $B$ is compact. In particular, $\pi(B)\subset U$ is compact.

Apply Lemma~\ref{lemma2.11} to $K:=\pi(B)\subset U$ and the chosen $v$ to obtain a~Delzant polytope~$P$ with
$
\pi(B)\subset \operatorname{Int} P\subset U
$
and with a~facet $F\subset \partial P$ whose outward primitive normal is~$v$.
Let~${(X_P,\omega_P,\mu_P)}$ be the associated compact symplectic toric manifold with moment polytope~$P$.
By the action--angle identification, $\mu_P^{-1}(\operatorname{Int}P)$ is symplectomorphic to $\operatorname{Int}P\times\bT^n$; we use this to view $B$ as a~subset of $X_P$.
Extend $J|_{B}$ arbitrarily to an $\omega_P$-tame almost complex structure~$\overline{J}$ on $X_P$.

Let $A:=\min_{F'}\omega_P(\beta_{F'})$, where $F'$ ranges over facets of $P$.
Since $F$ is one of the facets, we have $A\le \omega_P(\beta_F)$.
Moreover, because $P\subset U$ and the facet $F$ has outward normal $v$, the defining constant $c_P$ of that facet satisfies $c_P\ge c_{\partial U}(v)=c$, hence
\[
\omega_P(\beta_F)=\langle v,p\rangle-c_P \le \langle v,p\rangle-c = \delta_{\partial U}(p;v) \le \delta+\varepsilon.
\]
By Proposition~\ref{prop:toric-disc-Psi} applied to $(X_P,\omega_P,L_p)$ and $\overline{J}$, there exists a~stable $\overline{J}$-holomorphic map with a~disc component $u_D$ and $\partial u_D\subset L_p$ satisfying
\[
E(u_D)=\int_D u_D^*\omega_P \le A \le \delta+\varepsilon.
\]
The disc component $u_D$ must meet the toric boundary
$
X_P\setminus \mu_P^{-1}(\operatorname{Int}P)=\mu_P^{-1}(\partial P)$.
Indeed, one can pick a~primitive of $\omega_P$ on $\mu_P^{-1}(\operatorname{Int}P)$ for which $L_p$ is exact, showing that any disc with boundary on $L_p$ with positive symplectic area must meet the toric boundary.

On the other hand, since $g_{\overline{J}}$ is $a$-bounded on $B=B_R(L_p)$, Lemma~\ref{lem:local-confinement} implies that $u_D(D)$ is contained in the ball of radius \smash{$\rho=\max\bigl(2CE(u_D)/r_0,\sqrt{CE(u_D)}\bigr)$} about $L_p$, with $r_0=r_0(a)$ and the universal constant $C$ from that lemma. Since $u_D$ meets $\mu_P^{-1}(\partial P)$ while, by construction, $B_R(L_p)$ does not, we must have $R\le\rho\le C(a)\max\bigl(\delta+\varepsilon,\sqrt{\delta+\varepsilon}\bigr)$.
Since $\varepsilon>0$ was arbitrary, this gives the claim.
\end{proof}

\begin{proof}[Proof of Proposition~\ref{prop:flux-distance}]
Apply Proposition~\ref{prop:flux-distance-strong} with $V = \bR^n\times\bT^n$. Since $J$ is $\omega$-compatible and $g_J$ is $a$-bounded globally on $\bR^n\times\bT^n$, it is $a$-bounded on $U\times\bT^n$. Moreover, $g_J$ is complete on $V$. The conclusion of Proposition~\ref{prop:flux-distance-strong} exactly yields the desired inequality.
\end{proof}

\begin{proof}[Proof of Theorem~\ref{tm:2}]
Suppose for contradiction that $U\neq \bR^n$. Since $U$ is open and convex, its boundary is nonempty, so $\delta_{\mathrm{flux}}(p;\partial U)$ is finite for any $p\in U$.
The hypothesis states that there exists a~geometrically bounded $\omega$-compatible almost complex structure $J$ on $U\times\bT^n$. This means $g_J$ is complete on $U\times\bT^n$ and is $a$-bounded for some $a>0$.
Apply Proposition~\ref{prop:flux-distance-strong} with~${V = U\times\bT^n}$. Since $V\setminus (U\times\bT^n) = \varnothing$, the distance $d_{g_J}(L_p, V\setminus (U\times\bT^n))$ is infinite. However, Proposition~\ref{prop:flux-distance-strong} provides a~finite upper bound in terms of $\delta_{\mathrm{flux}}(p;\partial U)$, which is a~contradiction. Hence $U=\bR^n$.
\end{proof}

I used AI-assisted tools for editing and improving the exposition (in particular for wording, structure, and catching typographical issues).
All mathematical content, results, and any remaining errors are my own responsibility.

\begin{thebibliography}{99}
\footnotesize\itemsep=0pt

\bibitem{AurouxAnticanonicalTDuality}
Auroux D., Mirror symmetry and {$T$}-duality in the complement of an
 anticanonical divisor, \textit{J.~G\"okova Geom. Topol. GGT} \textbf{1}
 (2007), 51--91, \href{http://arxiv.org/abs/0706.3207}{arXiv:0706.3207}.

\bibitem{ChoOhToric}
Cho C.-H., Oh Y.-G., Floer cohomology and disc instantons of {L}agrangian torus
 fibers in {F}ano toric manifolds,
 \href{https://doi.org/10.4310/AJM.2006.v10.n4.a10}{\textit{Asian~J. Math.}}
 \textbf{10} (2006), 773--814,
 \href{http://arxiv.org/abs/math.SG/0308225}{arXiv:math.SG/0308225}.

\bibitem{EntovGanorMembrez}
Entov M., Ganor Y., Membrez C., Lagrangian isotopies and symplectic function
 theory, \href{https://doi.org/10.4171/CMH/451}{\textit{Comment. Math. Helv.}}
 \textbf{93} (2018), 829--882,
 \href{http://arxiv.org/abs/1610.09631}{arXiv:1610.09631}.

\bibitem{FOOO_Toric1}
Fukaya K., Oh Y.-G., Ohta H., Ono K., Lagrangian {F}loer theory on compact toric
 manifolds.~{I}, \href{https://doi.org/10.1215/00127094-2009-062}{\textit{Duke
 Math.~J.}} \textbf{151} (2010), 23--174,
 \href{http://arxiv.org/abs/0802.1703}{arXiv:0802.1703}.

\bibitem{Gallot}
Gallot S., Hulin D., Lafontaine J., Riemannian geometry, 3rd ed., \textit{Universitext},
 \href{https://doi.org/10.1007/978-3-642-18855-8}{Springer}, Berlin, 2004.

\bibitem{GromanFloerOpen}
Groman Y., Floer theory and reduced cohomology on open manifolds,
 \href{https://doi.org/10.2140/gt.2023.27.1273}{\textit{Geom. Topol.}}
 \textbf{27} (2023), 1273--1390,
 \href{http://arxiv.org/abs/1510.04265}{arXiv:1510.04265}.

\bibitem{GromanWrappedSemiToricSYZ}
Groman Y., The wrapped {F}ukaya category for semi-toric~{C}alabi--{Y}au,
 \href{https://doi.org/10.4171/jems/1466}{\textit{J.~Eur. Math. Soc. (JEMS)}}
 \textbf{26} (2024), 2373--2439,
 \href{http://arxiv.org/abs/1805.03635}{arXiv:1805.03635}.

\bibitem{GromanVarolgunesCompleteEmbeddings}
Groman Y., Varolgunes U., Locality of relative symplectic cohomology for
 complete embeddings,
 \href{https://doi.org/10.1112/s0010437x23007492}{\textit{Compos. Math.}}
 \textbf{159} (2023), 2551--2637,
 \href{http://arxiv.org/abs/2110.08891}{arXiv:2110.08891}.

\bibitem{Gunther}
G{\"u}nther P., Einige S{\"a}tze {\"u}ber das {V}olumen von geod{\"a}tischen
 {K}ugeln, \textit{Ber. Verh. S{\"a}chs. Akad. Wiss. Leipzig Math.-Nat. Kl.}
 \textbf{104} (1960), 1--38.

\bibitem{KarshonTangExcision}
Karshon Y., Tang X., Symplectic excision,
 \href{http://arxiv.org/abs/2101.03534}{arXiv:2101.03534}.

\bibitem{McDuffSymplecticStructures}
McDuff D., Symplectic structures~-- a new approach to geometry, \textit{Notices
 Amer. Math. Soc.} \textbf{45} (1998), 952--960.

\bibitem{McDuffSalamonJHol}
McDuff D., Salamon D., {$J$}-holomorphic curves and symplectic topology, 2nd ed.,
 \textit{Amer. Math. Soc. Colloq. Publ.}, Vol.~52, American
 Mathematical Society, Providence, RI, 2012.

\bibitem{STVFluxShape}
Shelukhin E., Tonkonog D., Vianna R., Geometry of symplectic flux
 and~{L}agrangian torus fibrations,
 \href{https://doi.org/10.1112/topo.70002}{\textit{J.~Topol.}} \textbf{17}
 (2024), e70002, 56~pages,
 \href{http://arxiv.org/abs/1804.02044}{arXiv:1804.02044}.

\bibitem{Sikorav}
Sikorav J.-C., Some properties of holomorphic curves in almost complex
 manifolds, in Holomorphic {C}urves in {S}ymplectic {G}eometry, \textit{Progr.
 Math.}, Vol.~117,
 \href{https://doi.org/10.1007/978-3-0348-8508-9_6}{Birkh\"auser}, Basel,
 1994, 165--189.

\bibitem{StratmannParamRays}
Stratmann B., Removing parametrized rays symplectically,
 \href{https://doi.org/10.4310/jsg.2022.v20.n2.a4}{\textit{J.~Symplectic
 Geom.}} \textbf{20} (2022), 499--508,
 \href{http://arxiv.org/abs/2009.05465}{arXiv:2009.05465}.

\end{thebibliography}
\end{document}
